% concatenated from divcurlsphere_eng.tex
\documentclass[a4paper]{article}
\usepackage[T2A]{fontenc}
%\usepackage[cp1251]{inputenc}
\usepackage[english]{babel}
\usepackage{mathtools}
\usepackage{amsmath}
\usepackage{amsfonts,amssymb}
\usepackage{amsthm}
\usepackage{esint}

\usepackage{mathrsfs}

\usepackage{graphicx}

\textwidth=15cm
%\textheight=22cm


\hoffset 0mm
\voffset 0mm

\let\rom=\textup

\newcommand{\R}{\mathbb{R}}
\newcommand{\N}{\mathbb{N}}
\newcommand{\Cc}{\mathbb{C}}
\newcommand{\Z}{\mathbb{Z}}
\newcommand{\Ker}{{\rm Ker}}
\newcommand{\Img}{{\rm Img}}
\newcommand{\Span}{{\rm span}}
\newcommand{\rank}{{\rm rank}}
\newcommand{\real}{{\rm Re}}
\newcommand{\X}{\mathcal{X}}
\newcommand{\Uu}{\mathcal{U}}
\newcommand{\bx}{\mathbf{x}}
\newcommand{\bh}{\mathbf{h}}
\newcommand{\bff}{\mathbf{f}}
\newcommand{\dx}{\operatorname{\mathbf{dx}}}
\newcommand{\bu}{\mathbf{u}}
\newcommand{\ba}{\mathbf{a}}
\newcommand{\supp}{\operatorname{supp}}
\newcommand{\curl}{\operatorname{curl}}
\newcommand{\divv}{\operatorname{div}}
%\newcommand{\Ker}{\operatorname{Ker}}



\newcommand{\bv}{\mathbf{v}}
 \newcommand{\by}{\mathbf{y}} 
 \newcommand{\bz}{\mathbf{z}}
 \newcommand{\bn}{\mathbf{n}}

\newtheorem{theorem}{Theorem}[section]
\newtheorem{definition}{Definition}[section]
\newtheorem{remark*}{Remark}[]

\overfullrule5pt
%\mag1096
\mag1300

\title{On the Multi-Dimensional Divergence-Curl Problem and Its Connection with Pseudo-Harmonic Fields}
\author{A.V.\,Gorshkov}

\begin{document}
\maketitle
\abstract{This article addresses the solvability of the multi-dimensional divergence-curl problem with a no-slip boundary condition. A solvability criterion is derived as an orthogonality condition of the vorticity function to pseudo-harmonic fields. A countable family of such fields, sufficient for the solvability of the three-dimensional problem in the exterior of a sphere, is also presented.}

\section{Introduction}
This article studies the problem of finding a solenoidal vector field $\bv$ in a domain $\Omega$ with a prescribed curl function $\bff$
\begin{align}\label{mainsyst}
\begin{cases}
\curl \bv = \bff,\\
\divv \bv = 0,
\end{cases}
\end{align}
subject to the homogeneous boundary condition
\begin{align}\label{boundcond}
\bv|_{\partial \Omega} = 0.
\end{align}

This boundary value problem has applications in hydrodynamics, where the homogeneous Dirichlet condition --- known as the {\it no-slip condition} --- describes the interaction of a viscous medium with a solid body. The divergence-curl problem has been studied in various formulations with weaker boundary conditions. In \cite{KirchhartSchulz} it was studied without boundary conditions. The Cauchy problem for the three-dimensional divergence-curl system and $L_p$--$L_q$ estimates for the solution were obtained in \cite{GallayWayne}. Boundary value problems with prescribed tangential and normal boundary conditions were considered in \cite{Divcurl1}. The exterior Neumann problem was investigated in \cite{Divcurl2}, \cite{Divcurl3}. The magnetostatics problem with various boundary conditions was studied in variational form in \cite{BramblePasciak}. In the author's work \cite{G1}, the divergence-curl problem for incompressible flows with the no-slip condition was solved. The present paper extends that study to the multi-dimensional case.

As is well known, solvability of the two-dimensional divergence-curl problem requires orthogonality of $\bff$ to harmonic functions \cite{G1}\cite{Quartapelle2}. In this article we extend this condition to higher dimensions by replacing harmonic fields in the orthogonality condition with {\it pseudo-harmonic} ones.

The vector Laplacian on $k$-forms admits the following invariant representation:
$$
\Delta = d^*d + dd^*.
$$
Here $d$ is the exterior derivative and $d^*$ is the codifferential. The kernel of this operator consists of {\it harmonic fields}.

For scalar functions and 2-forms, the two-dimensional Laplacian has only one non-zero term in each differential--codifferential pair:
$\Delta = d^*d = \divv(\nabla)$, or $\Delta = dd^* = \curl(\nabla^\perp)$.

We will prove solvability of the divergence-curl problem for {\it pseudo-harmonic fields}, in which only one pair of differentials --- namely $dd^*$ --- vanishes. For 2-forms on planar domains, these pseudo-harmonic fields coincide with harmonic functions, making the orthogonality principle with respect to pseudo-harmonic fields uniform across all dimensions.

On manifolds without boundary, by the adjointness of $d$ and $d^*$, the condition $dd^* = 0$ is equivalent to $d^* = 0$. However, for boundary value problems the situation changes significantly: the weaker requirement of orthogonality of $\bff$ to the kernel of $dd^*$ turns out to be not only necessary but also sufficient.

An important question is also the construction of a countable family of pseudo-harmonic fields, orthogonality to which is equivalent to the homogeneous Dirichlet boundary condition. For planar domains, a countable set of orthogonality conditions to harmonic functions takes the form ($\bx=(x_1,x_2)$):
\begin{itemize}
\item inside a disk:
\begin{align*}
\int_{|\bx|<r_0} f(\bx)(x_1+ix_2)^k \dx = 0,~k\in \N \cup \{0\},
\end{align*}

\item in the exterior of a disk:
\begin{align*}
\int_{|\bx|>r_0} \frac {f(\bx)}{(x_1+ix_2)^k} \dx = 0,~k\in \N \cup \{0\},
\end{align*}

\item in an annulus:
\begin{align*}
\int_{r_0<|\bx|<r_1} \frac {f(\bx)}{(x_1+ix_2)^k} \dx = 0,~k\in \Z.
\end{align*}
\end{itemize}
In the general case, analogous conditions for two-dimensional domains can be obtained via the Riemann mapping (see \cite{G1} for details).

The article presents one example of a three-dimensional divergence-curl problem with a countable family of pseudo-harmonic fields $\{\ba_n\}_{n\in \N}$ satisfying
$$
d d^* \ba_n = \curl \curl \ba_n =0.
$$
The orthogonality condition $(\bff,\ba_n)_{L_2}=0$ will be sufficient for the solvability of the problem.

This is particularly relevant for the numerical solution of the divergence-curl boundary value problem with the homogeneous no-slip condition, where the constraint on the right-hand side is imposed not on an entire countable family but only on a finite subset thereof, making the boundary value of the solution approximately zero --- a condition known in hydrodynamics as the {\it partial-slip condition}.

\section{The Divergence-Curl Problem in Differential Forms}

The divergence-curl problem consists of finding a form $\omega$ on a Riemannian manifold $M$ satisfying the system
\begin{align} \label{divcurl}
\begin{cases}
d\omega = f,\\
d^*\omega = 0.
\end{cases}
\end{align}

Here $\omega$ is the unknown form in the space of smooth $k$-forms $\Omega^k(M)$, $d:\Omega^k(M) \to \Omega^{k+1}(M)$ is the exterior derivative, $d^*:\Omega^k(M) \to \Omega^{k-1}(M)$ is the codifferential defined by $d^*=(-1)^{n+nk+1} *d*$, $*:\Omega^k(M) \to \Omega^{n-k}(M)$ is the Hodge star operator.
The inner product on forms is defined by
$(\omega_1, \omega_2) = \int_M \omega_1 \wedge *\omega_2$.

The exterior derivative defines the cochain complex
$$
{\displaystyle \ldots \xrightarrow {} \Omega_{n-1}(M)\xrightarrow {d_{n}} \Omega_{n}(M)\xrightarrow {d _{n+1}} \Omega_{n+1}(M)\xrightarrow {} \ldots },
$$
while $d^*$ defines the chain complex
$$
{\displaystyle \ldots \xleftarrow {} \Omega_{n-1}(M)\xleftarrow {d^*_{n}} \Omega_{n}(M)\xleftarrow {d ^*_{n+1}} \Omega_{n+1}(M)\xleftarrow {} \ldots }.
$$

Since $d^2=0$, a necessary condition for solvability is $df=0$. On manifolds without boundary, the codifferential $d^*$ is adjoint to $d$:
\begin{align}\label{scalprod}
(d\omega_1, \omega_2) = (\omega_1, d^*\omega_2).
\end{align}
Consequently, $f$ must be orthogonal to all forms $\omega_1$ in the kernel of $d^*$:
\begin{align}\label{ort1}
(f,\omega_1) = (d\omega, \omega_1) = (\omega, d^*\omega_1) = 0~\forall \omega_1:~d^*\omega_1=0.
\end{align}

Suppose the geometry of $M$ is such that $d^*\omega=0$ implies the existence of a potential $a$ satisfying $\omega = d^*a$. Note that such a form $a$ may exist even for domains with non-trivial topology under appropriate boundary conditions.
Then from the divergence-curl system we obtain the equation
$$d d^* a= f,$$
whose solvability requires $f$ to be orthogonal to all forms in the kernel of $dd^*$:
\begin{align}\label{ort2}
( f,\omega_1) = 0~\forall \omega_1:~dd^*\omega_1=0.
\end{align}

\begin{definition} A form $\omega$ is called \emph{pseudo-harmonic} if $dd^*\omega=0$.
\end{definition}

On the other hand, $(dd^* \omega_1, \omega_2) = (\omega_1, dd^* \omega_2)$, and from (\ref{scalprod}) it follows that on manifolds without boundary the condition $dd^*\omega_1=0$ is equivalent to the condition (\ref{ort1}), so no additional orthogonality constraint arises.

For manifolds with boundary, however, condition (\ref{ort2}) becomes stronger than (\ref{ort1}). By Stokes' formula:
\begin{align*}
\int_M (dd^* a,  f) = 
\int_M ( a, dd^*  f) +
\int_{\partial M} (d^* a, * f)
- \int_{\partial M} (d^* f, *a).
\end{align*}

We supplement the system (\ref{divcurl}) with the homogeneous boundary condition
$$
\omega\big|_{\partial M} = 0.
$$
Then $d^* a =0$ on $\partial M$. Assuming the potential $a$ can be chosen so that $* a = 0$ on $\partial M$, orthogonality to pseudo-harmonic forms (\ref{ort2}) becomes a necessary condition for solvability.

\vskip0.5cm
We illustrate the orthogonality principle in concrete dimensions: a solenoidal two-dimen\-si\-onal vector field $\bv=(v_1,v_2)$ with homogeneous boundary conditions can be represented as the skew gradient
$$\bv = d^*\psi = \nabla^\perp \psi$$
of a stream function $\psi$, where $\nabla_\bx^\perp = (-\partial_{x_2}, \partial_{x_1})$.

The no-slip condition $\bv=0$ on $\partial \Omega$ implies boundary conditions on the stream function:
\begin{align}\label{streamboundarycond}
\psi(\bx)= \mathrm{const}, \quad \frac {\partial \psi (\bx)}{\partial \bn} =0,~\bx \in \partial \Omega.
\end{align}
Here and below $\bn$ denotes the outward unit normal to the boundary.

Multiplying the Poisson equation
$$
\Delta \psi = \curl \bv
$$
by a harmonic function $h$ and applying Green's formula gives
$$
\int _{\Omega } \bff h \,\dx = \int _{\Omega }h \,\Delta \psi \,\dx=\int _{\Omega }\psi\,\Delta h\,\dx+\int _{\partial \Omega }\left(h{\frac {\partial \psi}{\partial \bn }}-\psi{\frac {\partial h}{\partial \bn }}\right)\,dS = 0.
$$
Orthogonality of $\bff$ to harmonic functions is both necessary and sufficient for solvability of (\ref{mainsyst}), (\ref{boundcond}).

\vskip 0.5cm
For $n=2$ the complexes take the form:
\begin{align*}
&0 \xrightarrow {} \Omega_{0}(M)\xrightarrow {\nabla} \Omega_{1}(M)\xrightarrow {\curl} \Omega_{2}(M)\xrightarrow {} 0, \\
&0 \xleftarrow {} \Omega_{0}(M)\xleftarrow {-\divv} \Omega_{1}(M)\xleftarrow {\nabla^\perp} \Omega_{2}(M)\xleftarrow {} 0.
\end{align*}

The Laplacian on $k$-forms is $\Delta = d^*d + dd^*$. For scalar functions it has only one non-zero pair:
$\Delta = d^*d = \divv(\nabla)$, or $\Delta = dd^* = \curl(\nabla^\perp)$, and $dd^*=0$ is equivalent to $\Delta=0$.
Therefore in two-dimensional divergence-curl problems, orthogonality to pseudo-harmonic fields reduces to orthogonality to harmonic functions, as described above.

\vskip 0.5cm
For $n=3$ the complexes are:
\begin{align*}
&0 \xrightarrow {} \Omega_{0}(M)\xrightarrow {\nabla} \Omega_{1}(M)\xrightarrow {\curl} \Omega_{2}(M)\xrightarrow {\divv} \Omega_{3}(M)\xrightarrow {} 0,\\
&0 \xrightarrow {} \Omega_{0}(M)\xleftarrow {-\divv} \Omega_{1}(M)\xleftarrow {\curl} \Omega_{2}(M)\xleftarrow {-\nabla} \Omega_{3}(M)\xleftarrow {} 0.
\end{align*}
The vector Laplacian contains both differential--codifferential pairs:
$$\Delta = d^*d + dd^* = -\nabla \divv + \curl^2.$$

If the flux through each connected component is zero, then a solenoidal field $\bv$ has a vector potential $\bv = \curl \ba$. Moreover, if $(\bv, \bn)=0$, the potential $\ba$ can be chosen satisfying $\ba \times \bn = 0$. Then $dd^* = \curl^2$ and condition (\ref{ort2}) becomes $\bff \perp \Ker \left ( \curl^2 \right )$.

\begin{definition} A covector field $\bv$ is called \emph{pseudo-harmonic} if
$$dd^* \bv = \curl^2 \bv = 0.$$
\end{definition}

\begin{remark*}
For dimensions greater than 2, forms lying in the kernels of both $d$ and $d^*$ are also subject to the orthogonality condition. On manifolds without boundary, harmonic fields lie in the kernels of both operators. However, orthogonality to harmonic fields is not sufficient for solvability of the boundary value problem (\ref{mainsyst}), (\ref{boundcond}). The broader class of pseudo-harmonic fields provides the correct class of forms appearing in the necessary and sufficient conditions for solvability.
\end{remark*}

\section{Solvability of the Divergence-Curl Problem}

Let $\Omega$ be a bounded Lipschitz domain, $\mathcal D(\Omega)$ the space of smooth compactly supported functions in $\Omega$. Define
$$
\mathcal{V}=\{ \bv \in \mathcal D(\Omega),~\divv \bv=0 \},
$$
and let $\mathcal H$ be its closure in $L_2(\Omega)$. Since $\divv \bv \in L_2(\Omega)$, the normal trace $(\bv,\bn)\in H^{-1/2}(\partial \Omega)$ is well defined.
We also introduce the spaces
\begin{align*}
&H(\curl,\Omega) = \{\ba\in L_2(\Omega),~\curl \ba\in L_2(\Omega)\},\\
&H(\curl^2,\Omega) = \{\ba \in L_2(\Omega),~\curl \ba\in L_2(\Omega),~\curl^2 \ba\in L_2(\Omega)\}.
\end{align*}
For $\bv \in H(\curl, \Omega)$, the trace $\bv \times \bn \in H^{-1/2}(\partial \Omega)$ is defined; for $\bv \in H(\curl^2, \Omega)$, so is $\curl \bv \times \bn \in H^{-1/2}(\partial \Omega)$. We then define
$$
H_0(\curl^2,\Omega) = \{\ba\in H(\curl^2,\Omega),~( \ba \times \bn)=0,~(\curl \ba \times \bn)=0\}
$$
as the closure of $\mathcal D(\Omega)$ under the norm
$$
\|\ba \|_{H_0(\curl^2,\Omega)}^2 = \|\ba \|_{H^1(\Omega)}^2 + \| \curl \ba \|_{L_2(\Omega)}^2+ \| \curl^2 \ba \|_{L_2(\Omega)}^2.
$$

On $L_2(\Omega)$ define the norm
$$
\|\bff\|_{H^*(\Omega)} = \sup_{\ba \in H_0(\curl^2),\ba\neq 0} \frac{\int_{\Omega} (\bff(\bx), \ba) \dx}{\|\ba\|_{H_0(\curl^2)}}.
$$
Let $H^*(\Omega)$ denote the completion of $L_2$ under this norm.

Before stating the solvability theorem, we give a definition of a solution. The boundary condition (\ref{boundcond}) implies that on each connected component $\partial \Omega_i$ of $\partial \Omega$, the sufficient condition for the existence of a vector potential holds:
$$
\int_{\partial \Omega_i} (\bv, \bn) \dx = 0.
$$
By the Weyl decomposition \cite{Monk}\cite{Temam}, any $\bv \in L_2$ then has a vector potential $\ba$ with $\bv = \curl \ba$, and this potential can be chosen satisfying $(\ba \times \bn)=0$ (see \cite{BramblePasciak}).

The problem then reduces to:
\begin{align}
&\curl^2 \ba = \bff, \label{maineq1} \\
&\curl \ba  \times \bn |_{\partial \Omega} = 0,\label{mainsystcond1}\\
&\ba  \times \bn |_{\partial \Omega} = 0.\label{mainsystcond2}
\end{align}

Condition (\ref{mainsystcond1}) means orthogonality of the vector field $\bv$ to the boundary, and (\ref{mainsystcond2}) means tangentiality. Together they ensure (\ref{boundcond}).

\begin{definition} A \emph{solution} $\bv$ of problem (\ref{mainsyst}), (\ref{boundcond}) is the curl of a function $\ba\in H(\curl^2, \Omega)$ satisfying (\ref{maineq1})--(\ref{mainsystcond2}).
\end{definition}

The problem is also overdetermined, as was the case in two dimensions for the Poisson equation with boundary conditions (\ref{streamboundarycond}). In this section we prove the following theorem on unique solvability:
\begin{theorem}\label{mainTh}
Let $\Omega$ be a bounded Lipschitz domain and $\bff\in L_2(\Omega)$. The problem (\ref{mainsyst}), (\ref{boundcond}) has a unique solution if and only if:
\begin{enumerate}
\item $\divv \bff = 0,$
\item $\int_{\Omega} (\bff,\bh)\dx=0~\forall \bh \in \Ker \left ( \curl^2 \right ).$
\end{enumerate}
\end{theorem}

\begin{proof}
Solenoidality of $\bff$ follows from properties of the exterior derivative. Necessity of the orthogonality condition follows from the definition of the solution and the generalized Stokes formula. Multiplying (\ref{maineq1}) by $\bh$ with $\curl^2 \bh =0$ and integrating by parts, all boundary terms vanish due to the boundary conditions, giving
\begin{align*}
\int_M (\curl^2 \ba, \bh) \dx = 
\int_M (\ba, \curl^2 \bh) \dx +
\int_{\partial M} (\curl \ba \times \bn, \bh)\,ds - 
\int_{\partial M} (\ba \times \bn, \curl \bh) = 0
\end{align*}
for all $\bh \in \Ker \left ( \curl^2 \right )$.

For sufficiency: the operator $\curl^2: H_0(\curl^2,\Omega)\to L_2(\Omega)$ has an adjoint $\curl^{2,*}: L_2(\Omega) \to H^*(\Omega)$. By the boundary conditions, $\curl^2$ and $\curl^{2,*}$ agree on smooth functions, and $\Img \curl^2 = \left ( \Ker \curl^{2,*} \right )^\perp$. Since $\curl^{2,*}$ is an epimorphism, for any $\bff \in L_2$ orthogonal to $\Ker \curl^{2,*}$, a solution $\ba \in H_0(\curl^2,\Omega)$ exists.

Uniqueness follows from Stokes' formula. Let $\ba$ be a solution with $\bff=0$. Testing against $\ba$:
\begin{align*}
\int_\Omega (\curl^2 \ba,\ba) \dx= 
\int_\Omega (\curl \ba, \curl \ba) \dx + 
\int_{\partial \Omega} (\curl \ba \times \bn, \ba) \dx =
\int_\Omega (\curl \ba, \curl \ba) \dx = \\ \|\curl \ba\|_{L_2} = 0.
\end{align*}
Hence $\bv = \curl \ba \equiv 0$.  
\end{proof}

\begin{remark*}
Since $\bff \in \mathcal H$, the boundary condition (\ref{boundcond}) additionally implies $(\bff,\bn)=0$. Indeed, if $(\bff,\bn) \neq 0$ on $\partial \Omega$, then the tangential part of $\bv$ on the boundary is non-zero. This condition does not appear explicitly in Theorem \ref{mainTh} because condition~2 already implies orthogonality to gradient fields, and combined with solenoidality, the Gauss--Ostrogradsky formula gives vanishing normal component of $\bff$.
\end{remark*}

\section{The Divergence-Curl Problem in \\Spherical Domains}

We now apply the general theory to a three-dimensional divergence-curl problem in the exterior of a sphere. We construct a countable family of pseudo-harmonic fields $\ba_n \in \Ker(\curl^2)$ that determine its solvability. These fields will also be solenoidal, hence harmonic ($\Delta \ba_n =0$), while $\curl \ba_n$ remains non-zero.

In the domain $B_{r_0}=\left \{\bx \in \R^3 ~\big | ~ |\bx|>r_0 \right \}$ for some fixed $r_0>0$, consider the exterior divergence-curl problem (\ref{mainsyst}) with no-slip boundary conditions on the sphere and at infinity:
\begin{align}
&\bv \big |_{|\bx|=r_0} = 0, \label{boundcondsphere} \\
&\bv \to 0,~|\bx|\to \infty. \label{boundcondinfsphere}
\end{align}

We seek the solution as an expansion in vector spherical harmonics:
\begin{align*}
\bv =\sum _{\ell =0}^{\infty }\sum _{m=-\ell }^{\ell }\left(V_{\ell m}^{r}(r)\mathbf {Y} _{\ell m}+V_{\ell m}^{(1)}(r)\mathbf {\Psi}_{\ell m}+V_{\ell m}^{(2)}(r)\mathbf {\Phi}_{\ell m}\right).
\end{align*}

The vector spherical harmonics are defined by
\begin{align*}
&\mathbf {Y} _{\ell m}=Y_{\ell }^{m}{\hat {\mathbf {r} }},\\
&\mathbf {\Psi } _{\ell m}=r\nabla Y_{\ell }^{m},\\
&\mathbf {\Phi } _{\ell m}=\mathbf {r} \times \nabla Y_{\ell }^{m}
\end{align*}
through the scalar spherical harmonic $Y_{\ell }^{m}$:
$$
Y_{\ell }^{m}(\theta ,\varphi )={\sqrt {{\frac {(2\ell +1)}{4\pi }}{\frac {(\ell -m)!}{(\ell +m)!}}}}\,P_{\ell }^{m}(\cos {\theta })\,e^{im\varphi }.
$$

The divergence and curl are expressed as follows:
\begin{align*}
&\nabla \cdot \bv =\sum _{\ell =0}^{\infty }\sum _{m=-\ell }^{\ell }\left({\frac {dV_{\ell m}^{r}}{dr}}+{\frac {2}{r}}V_{\ell m}^{r}-{\frac {\ell (\ell +1)}{r}}V_{\ell m}^{(1)}\right)\mathbf Y_{\ell m},
\\
&\nabla \times \bv = \sum _{\ell =0}^{\infty }\sum _{m=-\ell }^{\ell }\Bigg(-{\frac {\ell (\ell +1)}{r}}V_{\ell m}^{(2)}\mathbf {Y} _{\ell m}-\left({\frac {dV_{\ell m}^{(2)}}{dr}}+{\frac {1}{r}}V_{\ell m}^{(2)}\right)\mathbf {\Psi } _{\ell m}+\\
&\left(-{\frac {1}{r}}V_{\ell m}^{r}+ {\frac {dV_{\ell m}^{(1)}}{dr}}+{\frac {1}{r}}V_{\ell m}^{(1)}\right)\mathbf {\Phi } _{\ell m}\Bigg).
\end{align*}

Denoting the spherical components of $\bff$ by $(f_{\ell m}^r, f_{\ell m}^{(1)}, f_{\ell m}^{(2)})$, the system (\ref{mainsyst}) in terms of harmonic coefficients reads:
\begin{align}
&-{\frac {\ell (\ell +1)}{r}}V_{\ell m}^{(2)} = f_{\ell m}^r, \label{mainsyst_sphere_eq1}\\
&-{\frac {dV_{\ell m}^{(2)}}{dr}}-{\frac {1}{r}}V_{\ell m}^{(2)} = f_{\ell m}^{(1)}, \label{mainsyst_sphere_eq2}\\
&-{\frac {1}{r}}V_{\ell m}^{r}+ {\frac {dV_{\ell m}^{(1)}}{dr}}+{\frac {1}{r}}V_{\ell m}^{(1)} = f_{\ell m}^{(2)},\label{mainsyst_sphere_eq3}\\
&\frac {dV_{\ell m}^{r}}{dr}+{\frac {2}{r}}V_{\ell m}^{r}-{\frac {\ell (\ell +1)}{r}}V_{\ell m}^{(1)} = 0.\label{mainsyst_sphere_eq4}
\end{align}

From the first equation we obtain $f_{\ell m}^r(r_0) = 0$, i.e., $(\bff,\bn)=0$ on the boundary --- a condition that holds for arbitrary domains (see the remark in the previous section).

For spherical domains, an additional boundary condition on a second component of $\bff$ can be derived. The solenoidality condition gives:
\begin{align}\label{divf}
\frac {df_{\ell m}^{r}}{dr}+{\frac {2}{r}}f_{\ell m}^{r}-{\frac {\ell (\ell +1)}{r}}f_{\ell m}^{(1)} = 0,
\end{align}
from which a further boundary condition follows:
\begin{align*}
r_0 \frac {df_{\ell m}^r(r_0)}{dr} \Big |_{r=r_0} - \ell(\ell+1)f_{\ell m}^{(1)}(r_0) = 0.
\end{align*}

The system (\ref{mainsyst_sphere_eq1})--(\ref{mainsyst_sphere_eq4}), supplemented by solenoidality (\ref{divf}), is degenerate: equation (\ref{mainsyst_sphere_eq2}) follows from (\ref{mainsyst_sphere_eq1}) and (\ref{divf}), while the last two equations (\ref{mainsyst_sphere_eq3}), (\ref{mainsyst_sphere_eq4}) yield integral constraints on $f_{\ell m}^{(2)}$. Since $\mathbf{\Psi}_{\ell m} = \mathbf{\Phi}_{\ell m} = 0$ for $\ell=0$, we proceed with $\ell=1,\ldots,\infty$, $m\in\Z \cap [-\ell,\ell]$.

The solution basis of the system
\begin{align*}
&\frac {dV_{\ell m}^{r}}{dr}+{\frac {2}{r}}V_{\ell m}^{r}-{\frac {\ell (\ell +1)}{r}}V_{\ell m}^{(1)} = 0, \\
&-{\frac {1}{r}}V_{\ell m}^{r}+ {\frac {dV_{\ell m}^{(1)}}{dr}}+{\frac {1}{r}}V_{\ell m}^{(1)} = f_{\ell m}^{(2)}
\end{align*}
consists of two vector functions
\begin{align*}
\begin{bmatrix}-\frac{r^{-2-\ell}}{\ell+1} \\r^{-2-\ell}\\\end{bmatrix},
\quad
\begin{bmatrix}\frac{r^{\ell-1}}\ell \\r^{\ell-1}\\\end{bmatrix}.
\end{align*}

The system can be solved explicitly with the boundary condition at infinity (\ref{boundcondinfsphere}):
\begin{align}
&V_{\ell m}^{r} = \frac{\ell (\ell +1)}{r^{2+\ell}}\int_{r_0}^r
s^{2+\ell}  f_{\ell m}^{(2)}(s)\,ds +
\frac{\ell (\ell +1)}{r^{1-\ell}}\int_{r}^\infty
s^{1-\ell}  f_{\ell m}^{(2)}(s)\,ds,  \label{sol_sphere_1}\\
&V_{\ell m}^{(1)} = - \frac{\ell}{r^{2+\ell}}\int_{r_0}^r
s^{2+\ell}  f_{\ell m}^{(2)}(s)\,ds +
\frac{(\ell +1)}{r^{1-\ell}}\int_{r}^\infty
s^{1-\ell}  f_{\ell m}^{(2)}(s)\,ds. \label{sol_sphere_2}
\end{align}

The integral solvability condition in terms of spherical coefficients is:
\begin{align}\label{ortcondspherepolar}
\int_{r_0}^\infty
s^{1-\ell}  f_{\ell m}^{(2)}(s)\,ds = 0.
\end{align}

Returning to the original vector function $\bff$ via
\begin{align*}
f_{\ell m}^{(2)}(r) = \frac1{\ell (\ell +1)} \iint_{|\bx|=1} \bff(r\bx) \cdot \overline{\mathbf {\Phi}_{\ell m}}(\bx)\,dS,
\end{align*}
we obtain, for $\ell=1,\ldots,\infty$, $m\in\Z \cap [-\ell,\ell]$:
\begin{align}\label{ortcondsphere}
\iiint_{B_{r_0}} \frac{\bff(\bx) \cdot \overline{\mathbf {\Phi}_{\ell m}}\left (\bx/|\bx| \right )}
{|\bx|^{1 + \ell}} \dx =0.
\end{align}

Thus solvability requires $\bff(\bx)$ to be orthogonal to the family $\mathbf{\Phi}_{\ell m}/r^{\ell+1}$. We verify that these fields are indeed pseudo-harmonic, i.e.,
$$\curl^2 \frac{\mathbf \Phi_{\ell m}}{r^{\ell+1}} =0.$$

Using the curl identities for vector spherical harmonics (where $g(r)$ is a differentiable function):
\begin{align*}
\nabla \times \left(g(r)\mathbf {Y} _{\ell m}\right)&=-{\frac {1}{r}}g \mathbf {\Phi } _{\ell m},\\
\nabla \times \left(g(r) \mathbf {\Psi } _{\ell m}\right)&=\left({\frac {dg}{dr}}+{\frac {1}{r}}g\right)\mathbf {\Phi } _{\ell m},\\
\nabla \times \left(g(r)\mathbf {\Phi } _{\ell m}\right)&=-{\frac {\ell (\ell +1)}{r}}g\mathbf {Y} _{\ell m}-\left({\frac {dg}{dr}}+{\frac {1}{r}}g\right)\mathbf {\Psi } _{\ell m},
\end{align*}
we compute:
\begin{align*}
&\nabla \times \frac{\mathbf \Phi_{\ell m}}{r^{\ell+1}} = - \frac{\ell(\ell+1)}{r^{\ell+2}}\mathbf Y_{\ell m} + \frac \ell{r^{\ell+2}} \mathbf \Psi_{\ell m},\\
&\nabla \times \frac{\mathbf Y_{\ell m}}{r^{\ell+2}} = -\frac 1{r^{\ell+3}} \mathbf \Phi_{\ell m},\\
&\nabla \times \frac{\mathbf \Psi_{\ell m}}{r^{\ell+2}} =
-\frac {\ell+1}{r^{\ell+3}} \mathbf \Phi_{\ell m}.
\end{align*}

Hence:
$$
\nabla \times \nabla \times \frac{\mathbf \Phi_{\ell m}}{r^{\ell+1}} = -\frac{\ell(\ell+1) - \ell(\ell+1)}{r^{\ell+3}}\mathbf \Phi_{\ell m}  = 0.
$$

These fields are also divergence-free: since the radial component of $\mathbf{\Phi}_{\ell m}$ is zero, for any function $g(r)$
\begin{align*}
\nabla \cdot \left(g(r){\mathbf\Phi } _{\ell m}\right)&=0.
\end{align*}
Therefore
$$
\Delta \frac{\mathbf \Phi_{\ell m}}{r^{\ell+1}} = 0,
$$
and the fields $\frac{\mathbf \Phi_{\ell m}}{r^{\ell+1}}$ are harmonic --- yet their curl is non-zero.

\begin{theorem}
Let $\bff \in L_p(B_{r_0})$ with $p>1$, satisfying $\divv \bff = 0$, $(\bff,\bn)\big|_{|\bx|=r_0} = 0$, and the orthogonality condition (\ref{ortcondsphere}) for $\ell=1,\ldots,\infty$, $m\in\Z \cap [-\ell,\ell]$. Then there exists a unique solution of the divergence-curl problem (\ref{mainsyst}), (\ref{boundcond}) with boundary conditions (\ref{boundcondsphere}), (\ref{boundcondinfsphere}), given explicitly by the Biot--Savart-Laplace formula
\begin{align} \label{BiotSavar3d}
\bv(\bx) = -\frac{1}{4\pi} \iiint_{B_{r_0}} \frac{(\bx-\by)\times \bff(\by)}{|\bx-\by|^3} \,d\mathbf{y},
\end{align}
with estimates
\begin{align}
&\|\bv(\cdot)\|_{L_q} \leq C   \| \bff \|_{L_p},\quad\frac 1q = \frac 1p - \frac 13,\quad 1<p<q<\infty, \label{est1} \\
&\|\nabla \bv(\cdot)\|_{L_p} \leq C \| \bff \|_{L_p},\quad p>1. \label{est2}
\end{align}
\end{theorem}

\begin{proof}
The solution is given by the explicit formulas (\ref{sol_sphere_1}), (\ref{sol_sphere_2}). The boundary conditions (\ref{boundcondsphere}), (\ref{boundcondinfsphere}) are satisfied by construction. 
We extend the solution $\bv(\by)$ by zero inside the sphere. Then formula (\ref{BiotSavar3d}) becomes the Biot--Savart--Laplace law, which gives the explicit solution of the Cauchy problem for system (\ref{mainsyst}), (\ref{boundcondinfsphere}) (see \cite{GallayWayne}). We now prove that the solution extended by zero is also a solution of the exterior problem. In general, a solution of an elliptic system extended by zero in this way need not be a solution of the Cauchy problem. In spherical harmonics, the system (\ref{mainsyst}) reduces to a system of one-dimensional equations (\ref{mainsyst_sphere_eq1})--(\ref{mainsyst_sphere_eq4}) in the class of functions that are absolutely continuous on $(r_0,\infty)$ with derivative in $L_p$ for some $p>1$. Since the solution of this system vanishes at $r=r_0$ when the conditions of the theorem are satisfied, the solution extended by zero to $[0,\infty)$
 is also absolutely continuous, with derivative remaining in the same class $L_p$. The solution of the Cauchy problem (\ref{mainsyst}), (\ref{boundcondinfsphere}) is unique and is given by formula (\ref{BiotSavar3d}).


%Extending $\bv$ by zero inside the sphere, formula (\ref{BiotSavar3d}) becomes the Biot--Savart--Laplace law, giving the unique solution of the Cauchy problem for (\ref{mainsyst}), (\ref{boundcondinfsphere}) (see \cite{GallayWayne}). Since the solution vanishes at $r=r_0$ by the orthogonality conditions, the zero extension is absolutely continuous with derivative in $L_p$, and remains a solution.

The integral in (\ref{BiotSavar3d}) is a Riesz potential, so the first estimate (\ref{est1}) follows from the Hardy--Littlewood--Sobolev inequality (see [8, Chapter~V, \S~1]). The function $\nabla\frac{\bx}{|\bx|^3}$ is the kernel of a Calder\'{o}n--Zygmund singular integral operator, so the second estimate follows from the corresponding $L_p$ boundedness result (see Theorem~II.3 in \cite{Calderon}).
\end{proof}

\begin{remark*}
Unlike in a bounded domain, in an unbounded domain the orthogonality condition does not include orthogonality to constants: the requirement
$$\iiint_{B_{r_0}}\bff \,d\mathbf{x}=0$$
is not necessary. In two dimensions, for $\Omega \subset \R^2$ with the no-slip condition (\ref{boundcond}), Stokes' theorem gives a non-trivial circulation at infinity:
\begin{align*}
\lim_{R\to\infty}\oint_{\vert\bx\vert=R} \bv \cdot d\mathbf{l} = \iint_{\Omega}\bff \,d\mathbf{x} \neq 0.
\end{align*}
The $L_2$-norm of such solutions is infinite. In three dimensions:
\begin{align*}
\lim_{R\to\infty}\oiint_{\vert\bx\vert=R} \bv \times \bn \, dS = \iiint_{\Omega}\bff \,d\mathbf{x} \neq 0.
\end{align*}
A non-zero mean of the vorticity function produces circulatory flows at infinity while preserving (\ref{boundcondinfsphere}). Even so, from (\ref{est1}) $\bv$ may have finite $L_2$-norm if $\bff \in L_{6/5}$.
\end{remark*}

\vskip 0.5cm
Replacing condition (\ref{boundcondinfsphere}) with a prescribed uniform flow at infinity
\begin{align}
\bv \to \bv_\infty,~|\bx|\to \infty, \label{boundcondinfsphere2}
\end{align}
the solvability condition instead of (\ref{ortcondspherepolar}) becomes:
\begin{align*} 
\int_{r_0}^\infty
s^{1-\ell}  f_{\ell m}^{(2)}(s)ds = - \frac{r_0^{1-\ell}}{\ell (\ell +1)} v_{\infty, \ell m}^{r},
\end{align*}
or
\begin{align*}
\iiint_{B_{r_0}} \frac{\bff(\bx) \cdot \overline{\mathbf {\Phi}_{\ell m}}\left (\bx/|\bx| \right )}
{|\bx|^{1 + \ell}} \dx = - \frac{r_0^{1-\ell}}{\ell (\ell +1)} v_{\infty, \ell m}^{r},
\end{align*}
where $v_{\infty, \ell m}^{r}$ is the radial component of $\bv_\infty$ in the $\mathbf{Y}_{\ell m}$ direction. Since $v_{\infty, \ell m}^{r}$ vanishes except for isolated $\ell$, $m$, the orthogonality conditions largely remain intact.

The radial and two tangential components $\hat{\mathbf{r}}$, $\hat{\boldsymbol{\theta}}$, $\hat{\boldsymbol{\varphi}}$ of the vector field are related to Cartesian components $\{ \hat {e_k}\}$ by:
\begin{align*}
&\hat {\mathbf  r}=\sin \theta \cos \varphi {\boldsymbol {\hat {e_1 }}}+\sin \theta \sin \varphi {\boldsymbol {\hat {e_2 }}}+\cos \theta {\boldsymbol {\hat {e_3}}},\\
&\hat {\mathbf \theta }=\cos \theta \cos \varphi {\boldsymbol {\hat {e_1 }}}+\cos \theta \sin \varphi {\boldsymbol {\hat {e_2 }}}-\sin \theta {\boldsymbol {\hat {e_3}}},\\
&\hat {\mathbf \varphi }=-\sin \varphi {\boldsymbol {\hat {e_1 }}}+\cos \varphi {\boldsymbol {\hat {e_2 }}}.
\end{align*}

The vector $\bv_{\infty}$ with Cartesian coordinates $(0,0,1)$ has the spherical harmonic expansion:
$$
\bv_{\infty}=\cos \theta \,{\hat {\mathbf {r} }}-\sin \theta \,{\hat {\mathbf {\theta } }}=\mathbf {Y} _{10}+\mathbf{\Psi }_{10}.
$$
Hence $v_{\infty, \ell m}^{r}$ is non-zero only for $\ell=1$, $m=0$.





\section {Conclusion}
In dimensions greater than two vector fields from the kernel of $\curl^2$ (in dimension two - the kernel of $\Delta = \curl(\nabla^\perp)$) ensure the solvability of the divergence-curl problem with a homogeneous boundary condition. For certain spherical domains, it suffices to take a countable subset of vector spherical harmonics. These fields are solenoidal (but with non-zero curl), and consequently harmonic. In the exterior of a sphere, when the orthogonality condition is satisfied, the solution is given by the explicit Biot--Savart--Laplace formula.

\begin{thebibliography}{999}
\bibitem{KirchhartSchulz}Kirchhart M., Schulz E.: {\em Div-curl problems and -regular stream functions in 3D Lipschitz domains}. Math. Methods Appl. Sci., {\bf 45}, 1097--1117, (2021).

\bibitem{GallayWayne} Gallay Th., Wayne C.~E.: {\em Long-time asymptotics of the Navier-Stokes and vorticity equations on $\R^3$}. Philos. Trans. A Math. Phys. Eng. Sci.  {\bf 360}(1799), 2155--2188, (2002).

\bibitem{Divcurl1}Auchmuty, G., Alexander, J.~C.:{ \em{$L^2$ well-posedness of planar div-curl systems}}. {Arch. Ration. Mech. Anal.}, {\bf 160}, 91--134, (2001). 

\bibitem{Divcurl2}Neudert, M., von Wahl, W.: {\em{Asymptotic behaviour of the div-curl problem in exterior domains.}}
{Adv. Differ. Equ.}, {\bf 6}(11), 1347--1376, (2001).

\bibitem{Divcurl3} Delgado, B.~B., Mac\'ias-D\'iaz, J.~E.:  {\em{An exterior Neumann boundary-value problem for the div-curl system and applications}}. {Mathematics}, {\bf 9}, (2023). 

\bibitem{BramblePasciak} Bramble J.H., Pasciak J.E. {\em A new approximation technique for div-curl systems}. Mathematics of Computation. {\bf 73}(248), 1739--1762, (2004).

\bibitem{G1}Gorshkov, A.~V.: {\em {On the unique solvability of the div-curl problem in unbounded domains and energy estimates of solutions}}.  { Theor. Math. Phys.}, {\bf 221}, 1799--1812, (2024).  

\bibitem{Quartapelle2} Quartapelle, L., Valz-Gris, F.:  {\em{Projection conditions on the vorticity in viscous incompressible flows}}. { Internat. J. Numer. Methods Fluids}, {\bf 1}(2), 129--144, (1981).

\bibitem{Monk} P. Monk. {Finite Element Methods for Maxwell's Equations.} Clarendon Press, Oxford, 2003. 

\bibitem{Temam} R. Temam. {Navier-Stokes Equations Theory and Numerical Analysis.} North-Holland Publishing Company, New York, 1979.

\bibitem{Stein} E. Stein. {Singular integrals and differentiability properties of functions.} Princeton
University Press, Princeton, N.J., 1970. Princeton Mathematical Series, No. 30.

\bibitem{Calderon} Calder\'on, A. P., Zygmund, A.: {\em{On singular integrals.}} { Amer. J. Math.}, {\bf 78}(2), 289--309, (1956). 
\end{thebibliography}


\end{document}

